\chapter{HydroNodeNetzwerk}

\section*{Zielstellung}

Das HydroNodeNetzwerk bildet die kommunikationstechnische Grundlage des gesamten HydroNode-Systems.
Es verbindet physische Sensorstationen mit dem cloudbasierten Backend und der mobilen Anwendung.
Ziel des Netzwerks ist es, Messdaten von Hydroponik- und Umgebungssensoren zuverlässig, energieeffizient
und skalierbar zu übertragen, zu verarbeiten und für Nutzer verfügbar zu machen.

Das Netzwerk erfüllt dabei drei Kernaufgaben:

\begin{itemize}
  \item \textbf{Datenerfassung:} Periodisches Auslesen von Sensorwerten (Temperatur, Feuchte, pH-Wert u.\,a.)
        durch die eingebettete Hardware auf den Sensorstationen.
  \item \textbf{Datenübertragung:} Transport der Messdaten zum Backend über standardisierte Protokolle
        (HTTPS oder LoRaWAN), abhängig von Sensortyp und Infrastruktur.
  \item \textbf{Datenbereitstellung:} Persistierung, Anomalieerkennung und Ausgabe der Daten über
        eine REST-API an die iOS-Applikation \textit{HydroNode}.
\end{itemize}

Das Netzwerk ist so konzipiert, dass es heterogene Sensorhardware und unterschiedliche
Übertragungswege unter einer einheitlichen Abstraktion zusammenfasst.

\section*{Historische Entwicklung}

Das HydroNodeNetzwerk wurde ursprünglich für HTTPS-basierte Sensorik entwickelt.
Dabei kommunizieren Sensorstationen direkt über das Internet mit dem Backend-Endpunkt
\texttt{POST /api/webhook/sensor-value}.
Zur Vereinfachung der Entwicklung eigener Stationen wurde die
\textit{HydroNode Arduino Library}\footnote{\url{https://github.com/TexhFexLabs/HydroNode-Library}}
veröffentlicht, die HMAC-SHA256-Signierung, Zeitstempelgenerierung und Payload-Serialisierung
für Arduino-kompatible Mikrocontroller kapselt.

Im Zuge der Weiterentwicklung wurde das Netzwerk um LoRaWAN-Konnektivität erweitert.
Dabei wurde die bestehende Infrastruktur nicht ersetzt, sondern um einen zweiten
Übertragungskanal ergänzt.
Sensorstationen wie die \textit{HydroNodeStation01} können seitdem über das
Helium-LoRaWAN-Netzwerk Daten übertragen, ohne auf eine direkte Internetverbindung
angewiesen zu sein.
Die Integration ist abgeschlossen; offene Punkte betreffen ausschließlich das Fehlerhandling
für nicht verifizierbare Uplinks sowie die Dokumentation des Onboarding-Prozesses für
Drittanwender.

\newpage
\section*{Systemarchitektur}

Das HydroNodeNetzwerk gliedert sich in vier Schichten:

\begin{itemize}
  \item \textbf{Sensorschicht:} Eingebettete Hardware (Mikrocontroller, Sensoren, Funkmodul).
        Zwei Varianten: HTTPS-Stationen mit direkter Internetverbindung und LoRaWAN-Stationen
        mit Funkanbindung über Gateways.
  \item \textbf{Übertragungsschicht:} HTTPS (TLS 1.2+, HMAC-gesichert) oder LoRaWAN
        (AES-128, 868~MHz / 915~MHz je nach Region), Weiterleitung über Helium Network Server
        und ChirpStack Application Server zu Amazon SNS.
  \item \textbf{Backend-Schicht:} Quarkus-3-Anwendung (\textit{HydroNode Backend}).
        Empfang, Verifizierung, Kafka-basierte Verarbeitung und Persistierung in PostgreSQL.
  \item \textbf{Applikationsschicht:} iOS-App \textit{HydroNode}, die über die REST-API
        Messwerte, Stationsdaten und KI-Analysen abruft.
\end{itemize}

\begin{figure}[H]
  \centering
  \begin{tikzpicture}[
      font=\small,
      box/.style={draw, rounded corners=4pt, minimum height=11mm,
                  minimum width=40mm, align=center},
      sensor/.style={box, fill=blue!10, draw=blue!40!black},
      transport/.style={box, fill=orange!12, draw=orange!55!black},
      backend/.style={box, fill=green!10, draw=green!45!black},
      app/.style={box, fill=purple!10, draw=purple!45!black},
      flow/.style={-{Stealth[length=2.5mm]}, thick},
      layerlabel/.style={font=\footnotesize\itshape, text=gray!65!black, anchor=west}
  ]
    \node[sensor]    (https) at (-2.6,  0.0) {HTTPS-Station\\Arduino / ESP32};
    \node[sensor]    (lora)  at ( 2.6,  0.0) {LoRaWAN-Station\\HydroNodeStation01};

    \node[transport] (inet)  at (-2.6, -2.2) {HTTPS\\TLS\,1.2+ / HMAC};
    \node[transport] (hel)   at ( 2.6, -2.2) {Helium $\to$ ChirpStack\\$\to$ Amazon SNS};

    \node[backend, minimum width=84mm] (back) at (0, -4.4)
      {HydroNode Backend\\Kafka $\cdot$ PostgreSQL $\cdot$ Verifikation};

    \node[app, minimum width=62mm] (appn) at (0, -6.2)
      {iOS-App \textit{HydroNode} \enspace REST-API};

    \node[layerlabel] at (5.0,  0.0) {Sensorschicht};
    \node[layerlabel] at (5.0, -2.2) {Übertragungsschicht};
    \node[layerlabel] at (5.0, -4.4) {Backend-Schicht};
    \node[layerlabel] at (5.0, -6.2) {Applikationsschicht};

    \draw[flow, blue!60!black]   (https) -- (inet);
    \draw[flow, orange!60!black] (lora)  -- (hel);
    \draw[flow, blue!60!black]   (inet)  -- (back.north -| inet);
    \draw[flow, orange!60!black] (hel)   -- (back.north -| hel);
    \draw[flow, green!45!black]  (back)  -- (appn);
  \end{tikzpicture}
  \caption{Systemarchitektur des HydroNodeNetzwerks}
  \label{fig:hydronode-architektur}
\end{figure}

\newpage

\section*{Datenfluss}

\subsection*{HTTPS-Sensorpfad}

\begin{enumerate}
  \item Die Sensorstation misst periodisch Umgebungswerte und serialisiert diese als JSON-Payload.
  \item Der Payload wird mit dem geräteindividuellen Secret mittels HMAC-SHA256 signiert.
        Der Signaturwert wird im HTTP-Header \texttt{X-Signature} übertragen;
        ein Unix-Zeitstempel im Header \texttt{X-Timestamp} schützt vor Replay-Angriffen
        (Toleranz: $\pm20$~Minuten).
  \item Das Backend empfängt den Request unter \texttt{POST /api/webhook/sensor-value},
        publiziert den Rohpayload in das Kafka-Topic \texttt{hydronode.ingress.sensor} und
        antwortet sofort mit \texttt{202 Accepted}.
  \item Der \texttt{DeviceKeyVerifier}-Consumer liest den Ingress-Topic, prüft HMAC und
        Zeitstempel und publiziert bei Erfolg ein \texttt{SensorVerifiedEvent} nach
        \texttt{hydronode.verified.sensor}. Fehlerhafte Nachrichten landen im
        Dead Letter Queue-Topic \texttt{hydronode.invalid.sensor}.
  \item Der \texttt{SensorValueProcessor}-Consumer persistiert die Messwerte in PostgreSQL,
        aktualisiert die statistischen Baselines und übergibt neue Werte an den
        \texttt{AnomalyDetector}.
\end{enumerate}

\subsection*{LoRaWAN-Sensorpfad (Helium / ChirpStack~v4)}

\begin{enumerate}
  \item Die LoRaWAN-Station sendet einen verschlüsselten Uplink über das Helium-Funknetz.
        Das Helium Network Server leitet den Uplink an einen konfigurierten ChirpStack
        Application Server weiter.
  \item ChirpStack kodiert das Uplink-Event als JSON im ChirpStack-v4-Format und
        sendet es als SNS-Notification an das konfigurierte Amazon SNS Topic.
  \item Amazon SNS stellt die Nachricht per HTTP POST an den Backend-Endpunkt
        \texttt{POST /api/webhook/lorawan/sns} zu.
  \item Das Backend extrahiert \texttt{devEui}, \texttt{applicationId}, \texttt{fPort} sowie
        den Base64-kodierten Payload aus dem ChirpStack-Event.
        Über den \texttt{SensorLookupCache} wird die interne Sensor-UUID ermittelt.
  \item Der konfigurierte \texttt{LoRaFPortConfig}-Eintrag für den jeweiligen \texttt{fPort}
        definiert das Kanal-Layout (Byteanzahl, Vorzeichen, Skalierungsfaktor je Kanal).
        Der \texttt{LoRaFPortDecoder} dekodiert den Payload nach diesem Schema (Big-Endian).
  \item Die dekodierten Werte werden als \texttt{SensorVerifiedEvent} direkt nach
        \texttt{hydronode.verified.sensor} publiziert und vom \texttt{SensorValueProcessor}
        identisch zum HTTPS-Pfad weiterverarbeitet.
  \item Parallel aktualisiert der \texttt{StationSignalUpdateService} die Signalmetriken
        der Station (RSSI, SNR, Gateway-Koordinaten) anhand der \texttt{rxInfo}-Daten.
        Es wird das Gateway mit dem höchsten SNR-Wert ausgewählt (bei Gleichstand: höchster RSSI).
\end{enumerate}


\section*{Schnittstellen und Integrationspunkte}

\subsection*{HTTPS-Webhook}

Der Endpunkt \texttt{POST /api/webhook/sensor-value} empfängt Messdaten von HTTPS-fähigen
Stationen. Authentifizierung erfolgt über HMAC-SHA256; das Secret wird bei der Sensor-Aktivierung
generiert und ist nur dem Geräteeigentümer zugänglich (\texttt{GET /api/sensors/\{id\}/credentials}).

\subsection*{Amazon SNS / LoRaWAN}

Der Endpunkt \texttt{POST /api/webhook/lorawan/sns} nimmt SNS-Notifications entgegen.
Er behandelt drei SNS-Nachrichtentypen:

\begin{itemize}
  \item \texttt{SubscriptionConfirmation}: Das Backend bestätigt die SNS-Subscription automatisch,
        indem es die im Feld \texttt{SubscribeURL} enthaltene URL per GET-Request aufruft.
        Es ist kein manueller Eingriff erforderlich.
  \item \texttt{Notification}: Enthält im Feld \texttt{Message} das JSON-kodierte ChirpStack-v4-Uplink-Event.
  \item \texttt{UnsubscribeConfirmation}: Wird quittiert, ohne weitere Aktion.
\end{itemize}

\subsection*{Helium-Onboarding für Drittanwender}

Nutzer, die eine eigene LoRaWAN-Station über das Helium-Netzwerk anbinden möchten, folgen
diesen Schritten:

\begin{enumerate}
  \item In der Helium Console eine \textit{Integration} vom Typ \textit{HTTP} oder über einen
        verbundenen ChirpStack Application Server anlegen.
  \item In ChirpStack unter \textit{Applications $\rightarrow$ Integrations} eine
        \textit{AWS SNS}-Integration konfigurieren:
        AWS Region, SNS Topic ARN sowie IAM-Zugangsdaten mit \texttt{sns:Publish}-Berechtigung eintragen.
  \item Im AWS SNS-Topic eine Subscription vom Typ \textit{HTTPS} erstellen mit der Endpoint-URL:
\begin{verbatim}
https://hydronode.texhfexlabs.de/api/webhook/lorawan/sns
\end{verbatim}
        \textbf{Wichtig:} \textit{Raw Message Delivery} muss \textbf{deaktiviert} bleiben, da das Backend den vollständigen SNS-Envelope mit Metadaten benötigt.
  \item Nach Anlage der Subscription sendet AWS eine \texttt{SubscriptionConfirmation}-Nachricht
        an den Endpunkt. Das HydroNode-Backend bestätigt diese automatisch.
        Im AWS-Console-Status erscheint die Subscription anschließend als \textit{Confirmed}.
  \item Die DevEUI und ApplicationId des Geräts müssen im HydroNode-Backend dem Sensor zugeordnet
        sein (Felder \texttt{devEui} und \texttt{heliumApplicationId} in der \texttt{Sensor}-Entität).
  \item Für den genutzten \texttt{fPort} muss eine \texttt{LoRaFPortConfig} hinterlegt sein,
        die das Payload-Layout beschreibt (\texttt{POST /api/sensors/\{id\}/fport-configs}).
\end{enumerate}

\begin{figure}[H]
  \centering
  \resizebox{\linewidth}{!}{%
  \begin{tikzpicture}[
      font=\small,
      ext/.style={draw, rounded corners=4pt, minimum height=11mm,
                  minimum width=30mm, align=center, fill=gray!8},
      inn/.style={draw, rounded corners=3pt, fill=white,
                  minimum width=46mm, minimum height=9mm, align=center},
      flow/.style={-{Stealth[length=2.5mm]}, thick}
  ]
    \node[ext] (console) at (0, 0) {Helium Console};
    \node[ext, right=12mm of console] (chirp)  {ChirpStack\\App Server};
    \node[ext, right=12mm of chirp]   (awssns) {AWS SNS\\Topic};

    \node[inn, right=22mm of awssns]      (resource) {SnsWebhookResource};
    \node[inn, below=8mm of resource]     (lookups)  {SensorLookupCache\\FPortConfigCache};
    \node[inn, below=8mm of lookups]      (topic)    {Kafka: \texttt{verified.sensor}};

    \begin{pgfonlayer}{background}
      \node[draw, rounded corners=6pt, fill=green!8, draw=green!40!black,
            fit=(resource)(lookups)(topic), inner sep=9pt,
            label={[font=\small\bfseries, inner sep=5pt]above:HydroNode Backend}]
            (backbox) {};
    \end{pgfonlayer}

    \draw[flow] (console) -- (chirp);
    \draw[flow] (chirp)   -- (awssns);
    \draw[flow] (awssns)  -- (resource);
    \node[font=\footnotesize, above, align=center]
      at ($(awssns.east)!0.5!(backbox.west) + (0, 1.5)$)
      {HTTPS POST\\(auto-confirmed)};

    \draw[flow] (resource) -- (lookups);
    \draw[flow] (lookups)  -- (topic);

    \node[font=\footnotesize\itshape, below=5mm of backbox]
      {IAM-Policy: \texttt{sns:Publish}};
  \end{tikzpicture}}
  \caption{Integrationspfad: Helium / ChirpStack / Amazon SNS}
  \label{fig:helium-sns-integration}
\end{figure}

\subsection*{Kafka-interne Kommunikation}

Das Backend nutzt Apache Kafka als internen Nachrichtenbus:

\begin{table}[h]
  \centering
  \begin{tabular}{ll}
    \hline
    \textbf{Topic} & \textbf{Zweck} \\
    \hline
    \texttt{hydronode.ingress.sensor}    & Rohe HTTPS-Payloads (unverified) \\
    \texttt{hydronode.verified.sensor}   & Verifizierte Messwerte (beide Pfade) \\
    \texttt{hydronode.invalid.sensor}    & Dead Letter Queue (HMAC-Fehler, unbekannte Geräte) \\
    \texttt{hydronode.ingress.lorawan}   & Archiv aller LoRaWAN-Rohpayloads (fire-and-forget) \\
    \hline
  \end{tabular}
  \caption{Kafka-Topics im HydroNodeNetzwerk}
  \label{tab:kafka-topics}
\end{table}

\subsection*{REST-API / iOS-App}

Die iOS-App \textit{HydroNode} kommuniziert ausschließlich über die REST-API des Backends.
Relevante Endpunkte für das Netzwerkkonzept:

\begin{itemize}
  \item \texttt{GET /api/stations/\{id\}}: Stationsdaten inkl.\ Signalmetriken (RSSI, SNR,
        Gateway-Koordinaten, Zeitpunkt des letzten Uplinks).
  \item \texttt{GET /api/stations/\{id\}/measurements}: Aggregierte Messwerte aller
        Sensoren einer Station.
  \item \texttt{GET /api/stations/public}: Öffentlich sichtbare Stationen für
        kartenbasierte Darstellungen.
\end{itemize}

\newpage

\section*{Betrieb, Skalierung und Zuverlässigkeit}

\subsection*{Caching}

Zur Reduzierung von Datenbankabfragen im kritischen Empfangspfad setzt das Backend
auf drei In-Process-Caches mit einer TTL von jeweils 90~Minuten:

\begin{itemize}
  \item \texttt{DeviceKeyCache}: Abbildung Sensor-UUID auf Secret (HTTPS-Authentifizierung).
  \item \texttt{SensorLookupCache}: Abbildung DevEUI auf Sensor-UUID (LoRaWAN-Authentifizierung).
  \item \texttt{FPortConfigCache}: Abbildung (Sensor-UUID, fPort) auf Kanalstruktur
        (LoRaWAN-Dekodierung).
\end{itemize}

Cache-Invalidierung erfolgt nach jeder Konfigurationsänderung sowie manuell über die
Admin-API (\texttt{/api/admin/cache/*}).

\subsection*{Entkopplung durch Kafka}

Die Kafka-basierte Architektur entkoppelt Empfang und Verarbeitung von Nachrichten.
Ein kurzfristiger Ausfall des \texttt{SensorValueProcessor} führt nicht zu Datenverlust,
da Nachrichten im Topic erhalten bleiben bis sie konsumiert werden.
Der HTTPS-Webhook-Endpunkt antwortet dem Sender mit \texttt{202 Accepted}, sobald die
Nachricht im Ingress-Topic liegt, unabhängig von der weiteren Verarbeitungszeit.

\subsection*{Datenbankschema-Management}

Die Datenbankschemata werden automatisch durch Quarkus verwaltet.
In Entwicklungs- und Produktionsumgebung gilt die Strategie \texttt{update}:
neue Tabellen und Spalten werden ergänzt, bestehende Daten bleiben erhalten.
Manuelle SQL-Skripte sind nach Deployments nicht erforderlich, solange ausschließlich
additive Schemaänderungen vorgenommen werden.

\subsection*{Skalierungsüberlegungen}

Das System ist für den aktuellen Betriebsumfang (einzelne bis wenige Dutzend Stationen)
ausgelegt. Für eine horizontale Skalierung des Backends sind folgende Punkte relevant:

\begin{itemize}
  \item Die In-Process-Caches sind nicht über mehrere Backend-Instanzen hinweg synchronisiert.
        Bei mehreren Instanzen ist ein verteiltes Cache-Backend (z.\,B. Redis) erforderlich.
  \item Kafka-Consumer sind als Single-Consumer konfiguriert; eine Partitionierung
        der Topics ermöglicht parallele Verarbeitung durch mehrere Instanzen.
  \item Amazon SNS unterstützt nativ mehrere Subscriptions, sodass mehrere Backend-Instanzen
        parallel Uplinks empfangen können, sofern das Cache-Problem gelöst ist.
\end{itemize}

\newpage

\section*{Offene Punkte und Ausblick}

Folgende Aspekte sind in der aktuellen Implementierung noch nicht abschließend behandelt:

\begin{itemize}
  \item \textbf{Fehlerbehandlung nicht verifizierbarer Uplinks:} LoRaWAN-Uplinks, für die kein
        passender Sensor-Eintrag oder keine \texttt{fPortConfig} existiert, werden aktuell
        in den DLQ-Topic (\texttt{hydronode.invalid.sensor}) geschrieben. Eine automatische
        Benachrichtigung des Gerätebesitzers oder ein Retry-Mechanismus ist nicht implementiert.
  \item \textbf{Verteiltes Caching:} Bei Betrieb mehrerer Backend-Instanzen müssen die
        In-Process-Caches durch eine verteilte Lösung ersetzt werden.
  \item \textbf{Erweiterung der Arduino-Bibliothek:} Die HydroNode Arduino Library unterstützt
        aktuell ausschließlich den HTTPS-Pfad. Eine Erweiterung um LoRaWAN-spezifische
        Hilfsfunktionen (fPort-Payload-Kodierung) ist geplant.
  \item \textbf{Onboarding-Automatisierung:} Die manuelle Zuordnung von DevEUI und
        \texttt{fPortConfig} soll durch einen geführten Einrichtungsassistenten in der
        iOS-App vereinfacht werden.
\end{itemize}

\newpage
\section*{KI-gestützte Analyse in der HydroNode App}

  Die HydroNode iOS-App integriert Amazon Bedrock Nova Micro als KI-Komponente zur
  automatisierten Analyse von Sensoranomalien. Das Quarkus-Backend führt alle
  15~Minuten einen Analysezyklus durch. Erkennt die statistische Auswertung eine
  Auffälligkeit, etwa einen plötzlichen Wertesprung (\emph{SPIKE}), einen
  unerwarteten Einbruch (\emph{DROP}), einen schleichenden Drift
  (\emph{GRADUAL\_DRIFT}) oder eine Schwellenwertüberschreitung
  (\emph{THRESHOLD\_BREACH}), so wird die Anomalie serverseitig an Amazon Bedrock
  Nova Micro übergeben. Zusätzlich erzeugt das System alle 6~Stunden eine
  Routineanalyse aller aktiven Sensoren, auch ohne erkannte Anomalie, um
  schleichende Veränderungen frühzeitig zu erfassen.

  Dem Modell werden folgende Kontextdaten übergeben:

  \begin{itemize}
    \item 48~Stunden historische Messwerte: die letzten 3~Stunden als
          10-Minuten-Buckets, ältere Daten als Min/Max/Durchschnitt,
    \item referenzbasierte Schwellenwerte je Sensortyp
          (\emph{good} $|$ \emph{warn} $|$ \emph{crit}),
    \item die gelernte Baseline (\(\bar{x} \pm \sigma\), Stichprobenanzahl) aus
          einem rollierenden 72-Stunden-Fenster,
    \item Anomalietyp, Ist- und Erwartungswert sowie Abweichung in~$\sigma$,
    \item Sensorkanal mit Einheit (z.\,B.\ Temperatur in~°C,
          CO\textsubscript{2} in~ppm, pH-Wert),
    \item Erfassungszeitpunkt.
  \end{itemize}

  Der System-Prompt unterscheidet sich je nach Installation: Bei
  \emph{Outdoor}-Sensoren (Wetterstation) wird das Modell angewiesen,
  Ausreißer zunächst durch Wetterereignisse (Regen, Wind, Tagesgang) zu erklären
  und Alarme nur bei klar pathologischen Werten auszugeben.
  Bei \emph{Indoor}-Sensoren (Hydroponik, Grow-Tent) werden CO\textsubscript{2}-
  und VOC-Abfälle bei gleichzeitigem Temperaturwechsel als Lüftungsereignis
  gewertet (\emph{should\_alert\,=\,false}), während anhaltender
  CO\textsubscript{2}-Anstieg oder pH/EC-Drift als echtes Problem eingestuft wird.

  Das Modell liefert eine strukturierte JSON-Antwort zurück. Diese umfasst
  eine Erklärung der Ursache (\emph{explanation}), Konfidenz (\emph{confidence})
  und Umgebungsbewertung (\emph{environment\_score}) jeweils auf einer Skala
  von 0 bis 1, einen Luftqualitätsscore (\emph{air\_quality\_score}), die Erkennung
  eines Lüftungsereignisses (\emph{is\_ventilation\_event}), ein Alert-Flag
  (\emph{should\_alert}), einen Alert-Schweregrad
  (\emph{alert\_severity}: \texttt{NONE} $|$ \texttt{LOW} $|$ \texttt{MEDIUM}
  $|$ \texttt{HIGH} $|$ \texttt{CRITICAL}), eine Liste betroffener Sensortypen
  (\emph{affected\_sensors}) sowie eine konkrete Handlungsempfehlung
  (\emph{recommendation}).
  Die App stellt diese Ausgabe im \texttt{AnomalyDetailView} übersichtlich dar:
  Erklärungstext, qualitative Kennzahlen mit farbcodierten Badges sowie die
  priorisierte Maßnahme. Administratoren können über ein dediziertes
  Analytics-Dashboard den Token-Verbrauch und die geschätzten Kosten je
  Analysezeitraum einsehen.

      \section*{Datenvisualisierung und Benutzerinteraktion}
                                                            
  Die Detailansicht einer Sensorstation (\texttt{SensorDetailView}) stellt                                                                           
  Messzeitreihen als interaktiven Linien- und Flächengraphen dar, der auf
  dem Swift-Charts-Framework basiert. Für jeden Messkanal, darunter                                                                                 
  Temperatur, Luftfeuchte, pH-Wert und TDS, lässt sich ein Zeitfenster                                                                              
  zwischen~6~Stunden und~7~Tagen auswählen; ein benutzerdefinierter                                                                                  
  Datumsbereich ist ebenfalls möglich. Zur Performanzoptimierung aggregiert                                                                          
  die App Messwerte intervallabhängig (z.\,B.\ 20-Minuten-Mittel für                                                                                 
  24-Stunden-Ansichten). Erkannte Anomalien werden als farbige Markierungen                                                                          
  direkt im Graphen eingeblendet.                                                                                                                    
                                                                                                                                                     
  Für LoRaWAN-verbundene Stationen zeigt die                                                                                                         
  \texttt{StationSignalInfoView} die Funkqualitätsindikatoren RSSI                                                                                   
  (in~dBm) und SNR (in~dB) sowie die geografischen Koordinaten von Station                                                                           
  und Gateway. Eine eingebettete Kartenansicht visualisiert die                                                                                      
  Funkstrecke zwischen Sensorstation und nächstem Helium-Gateway.                                                                                    
                                                                                                                                                     
  Die Benutzerverwaltung umfasst das Teilen einer Station über                                                                                       
  \texttt{ShareSensorView}, die Zugriffsverwaltung per                                                                                               
  \texttt{SensorAccessListView} sowie die fPort-Konfiguration zur                                                                                    
  Steuerung des LoRaWAN-Dekodierungsverhaltens über                                                                                                  
  \texttt{FPortConfigsView}. Über das Kontextmenü der Detailansicht kann                                                                             
  der Nutzer außerdem Steuerbefehle in die Sensor-Queue einreihen                                                                                    
  (z.\,B.\ Pumpsteuerung) und die KI-Analyseansicht der jeweils                                                                                      
  ausgewählten Anomalie direkt aufrufen. 